\documentclass[12pt, titlepage]{article}

\usepackage{amsmath, mathtools}

\usepackage[round]{natbib}
\usepackage{amsfonts}
\usepackage{amssymb}
\usepackage{graphicx}
\usepackage{colortbl}
\usepackage{xr}
\usepackage{hyperref}
\usepackage{longtable}
\usepackage{xfrac}
\usepackage{tabularx}
\usepackage{float}
\usepackage{siunitx}
\usepackage{booktabs}
\usepackage{multirow}
\usepackage[section]{placeins}
\usepackage{caption}
\usepackage{fullpage}

\hypersetup{
bookmarks=true,     % show bookmarks bar?
colorlinks=true,       % false: boxed links; true: colored links
linkcolor=red,          % color of internal links (change box color with linkbordercolor)
citecolor=blue,      % color of links to bibliography
filecolor=magenta,  % color of file links
urlcolor=cyan          % color of external links
}

\usepackage{array}

\externaldocument{../../SRS/SRS}

\input{../../Comments.text}
\input{../../Common.text}

\newcounter{mnum}
\newcommand{\mthemnum}{M\themnum}
\newcommand{\mref}[1]{M\ref{#1}}

\begin{document}

\title{Module Interface Specification for \progname{}}

\author{\authname}

\date{\today}

\maketitle

\pagenumbering{roman}

\section{Revision History}

\begin{tabularx}{\textwidth}{p{3cm}p{2cm}X}
\toprule {\bf Date} & {\bf Version} & {\bf Notes}\\
\midrule
Date 1 & 1.0 & Notes\\
Date 2 & 1.1 & Notes\\
\bottomrule
\end{tabularx}

~\newpage

\section{Symbols, Abbreviations and Acronyms}

See SRS Documentation at \wss{give url}

\wss{Also add any additional symbols, abbreviations or acronyms}

\newpage

\tableofcontents

\newpage

\pagenumbering{arabic}

\section{Introduction}

The following document details the Module Interface Specifications for
Magic-8, a free-to-play, web-based 3D puzzle game built around an
8x8x8 magic cube. It states the purpose, users, and project goals: to provide an engaging,
navigable magic cube experience in which players complete tasks to practice addition and
subtraction in base 8, supported by an introductory tutorial, guided hints, online competitive
features, and task-based problem solving.

Complementary documents include the System Requirement Specifications
and Module Guide.  The full documentation and implementation can be
found at \url{https://github.com/Jacob-MB/magic-8}.

\section{Notation}


The structure of the MIS for modules comes from \citet{HoffmanAndStrooper1995},
with the addition that template modules have been adapted from
\cite{GhezziEtAl2003}.  The mathematical notation comes from Chapter 3 of
\citet{HoffmanAndStrooper1995}.  For instance, the symbol := is used for a
multiple assignment statement and conditional rules follow the form $(c_1
\Rightarrow r_1 | c_2 \Rightarrow r_2 | ... | c_n \Rightarrow r_n )$.

The following table summarizes the primitive data types used by \progname. 

\begin{center}
\renewcommand{\arraystretch}{1.2}
\noindent 
\begin{tabular}{l l p{7.5cm}} 
\toprule 
\textbf{Data Type} & \textbf{Notation} & \textbf{Description}\\ 
\midrule
character & char & a single symbol or digit\\
integer & $\mathbb{Z}$ & a number without a fractional component in (-$\infty$, $\infty$) \\
natural number & $\mathbb{N}$ & a number without a fractional component in [1, $\infty$) \\
real & $\mathbb{R}$ & any number in (-$\infty$, $\infty$)\\
\bottomrule
\end{tabular} 
\end{center}

\noindent
The specification of \progname \ uses some derived data types: sequences, strings, and
tuples. Sequences are lists filled with elements of the same data type. Strings
are sequences of characters. Tuples contain a list of values, potentially of
different types. In addition, \progname \ uses functions, which
are defined by the data types of their inputs and outputs. Local functions are
described by giving their type signature followed by their specification.

\section{Module Decomposition}

The following table is taken directly from the Module Guide document for this project.

\begin{table}[h!]
\centering
\begin{tabular}{p{0.3\textwidth} p{0.6\textwidth}}
\toprule
\textbf{Level 1} & \textbf{Level 2}\\
\midrule

Hardware-Hiding Module
&\refstepcounter{mnum} \mthemnum \label{mVH}: Virtual Hardware Module\\
\midrule

\multirow{11}{0.3\textwidth}{Behaviour-Hiding Module}
&\refstepcounter{mnum} \mthemnum \label{mUI}: User Interface Module\\
&\refstepcounter{mnum} \mthemnum \label{mRE}: Rendering Module\\
&\refstepcounter{mnum} \mthemnum \label{mN}: Navigation Module\\
&\refstepcounter{mnum} \mthemnum \label{mRC}: Round Control Module\\
&\refstepcounter{mnum} \mthemnum \label{mR}: Ruleset Module\\
&\refstepcounter{mnum} \mthemnum \label{mSC}: Scoring Module\\
&\refstepcounter{mnum} \mthemnum \label{mUM}: User Management Module\\
&\refstepcounter{mnum} \mthemnum \label{mL}: Leaderboard Module\\
&\refstepcounter{mnum} \mthemnum \label{mDC}: Daily Challenge Module\\
&\refstepcounter{mnum} \mthemnum \label{mT}: Tutorial Module\\
&\refstepcounter{mnum} \mthemnum \label{mSE}: Settings Module\\
\midrule

\multirow{8}{0.3\textwidth}{Software Decision Module}
&\refstepcounter{mnum} \mthemnum \label{mCR}: Cube Representation Module\\
&\refstepcounter{mnum} \mthemnum \label{mGDM}: Game Data Management Module\\
&\refstepcounter{mnum} \mthemnum \label{mRG}: Round Generation Module\\
&\refstepcounter{mnum} \mthemnum \label{mTG}: Line Geometry Module\\
&\refstepcounter{mnum} \mthemnum \label{mUDM}: User Data Management Module\\
&\refstepcounter{mnum} \mthemnum \label{mNBT}: Number Base Translation Module\\
&\refstepcounter{mnum} \mthemnum \label{mSCM}: Server Communication Module\\
&\refstepcounter{mnum} \mthemnum \label{mDCM}: Online Daily Challenge Manager Module\\
\bottomrule

\end{tabular}
\caption{Module Hierarchy}
\label{TblMH}
\end{table}


\newpage
~\newpage

\section{MIS of Virtual Hardware Module} \mref{mVH} 

\begin{description}
\item[Secrets:] The interface to the input and output devices.
\item[Services:]
\item[Implemented By:] pygame-ce, pygbag, and the web browser
\end{description}

\subsection{Module}

vhm

\subsection{Uses}

User Interface Module \mref{mUI}

Rendering Module \mref{mRE}

Navigation Module \mref{mN}

Round Control Module \mref{mRC}

Ruleset Module \mref{mR}

Scoring Module \mref{mSC}

User Management Module \mref{mUM}

Leaderboard Module \mref{mL}

Daily Challenge Module \mref{mDC}

Tutorial Module \mref{mT}

Settings Module \mref{mSE}


\subsection{Syntax}

\subsubsection{Exported Constants}



\subsubsection{Exported Access Programs}

\begin{center}
\begin{tabular}{p{2cm} p{4cm} p{4cm} p{2cm}}
\hline
\textbf{Name} & \textbf{In} & \textbf{Out} & \textbf{Exceptions} \\
\hline
main & - & - & - \\
\hline
\end{tabular}
\end{center}

\newpage

\section{MIS of User Interface Module} \mref{mUI} 

\begin{description}
\item[Secrets:] The mapping from user input to commands, and the presentation of the application state to the user, including the drawing of every screen.
\item[Services:] Reads device input and converts it into commands for \mref{mN}, converting typed answers from the display base with \mref{mNBT}. Draws each screen from the content provided by \mref{mN} and, for the game space, from the view produced by \mref{mRE}.
\end{description}

\subsection{Module}

ui

\subsection{Uses}

Rendering Module \mref{mRE}

Navigation Module \mref{mN}

Number Base Translation Module \mref{mNBT}


\subsection{Syntax}

\subsubsection{Exported Constants}

\begin{center}
\begin{tabular}{p{3cm} p{4cm} p{4cm} p{2cm}}
\hline
\textbf{Name} & \textbf{In} & \textbf{Out} & \textbf{Exceptions} \\
\hline


\hline
\end{tabular}
\end{center}

\subsubsection{Exported Access Programs}

\begin{center}
\begin{tabular}{p{4cm} p{4cm} p{4cm} p{4cm}}
\hline
\textbf{Name} & \textbf{In} & \textbf{Out} & \textbf{Exceptions} \\
\hline
load\_components & list & - & EmptyList, ComponentSyntaxError \\
verify\_components & Component Object & boolean & ValidationError, StateError, DependencyError, CacheError   \\
load\_component\_map & Coordinate\_Map\_List & - & EmptyList, MapSyntaxError \\
verify\_map & Coordinate Object & boolean & ValidationError, DependencyError, SizeError, CoordinateError \\

\hline
\end{tabular}
\end{center}

\newpage

\section{MIS of Rendering Module} \mref{mRE} 

\begin{description}
\item[Secrets:] The algorithms to translate the 3D space and game state into a 2D view, including layer isolation and highlighting.
\item[Services:] Produces the 2D view of the game space from the round state read from \mref{mRC}, for \mref{mUI} to display. It does not display the view itself.
\end{description}

\subsection{Module}

re

\subsection{Uses}

Round Control Module \mref{mRC}

User Interface Module \mref{mUI}



\subsection{Syntax}

\subsubsection{Exported Constants}

\begin{center}
\begin{tabular}{p{4cm} p{4cm} p{4cm} p{4cm}}
\hline
\textbf{Name} & \textbf{In} & \textbf{Out} & \textbf{Exceptions} \\
\hline


\hline
\end{tabular}
\end{center}

\subsubsection{Exported Access Programs}

\begin{center}
\begin{tabular}{p{4cm} p{4cm} p{4cm} p{4cm}}
\hline
\textbf{Name} & \textbf{In} & \textbf{Out} & \textbf{Exceptions} \\
\hline
space\_into\_2D & 3D Coordinates List & 2D Coordinates List & ConversionError, MissingCoordinateError \\
extract\_highlighted\_squares & Cell Array & Selection Array &  PointerOutOfBoundsError, TooManySelectionsError \\
layer\_highlights & $\mathbb{Z}$ & $\mathbb{Z}$ & ImpededLayerError \\


\hline
\end{tabular}
\end{center}

\newpage

\section{Navigation Module} \mref{mN} 

\begin{description}
\item[Secrets:] The structure of the application's screens and menus (including the settings screen) and the transitions between them, including those into and out of a game round.
\item[Services:] Tracks the current screen, carries out the commands passed in from \mref{mUI}, and provides the content of the current screen to \mref{mUI}.
\end{description}

\subsection{Module}

nav

\subsection{Uses}

User Interface Module \mref{mUI}



\subsection{Syntax}

\subsubsection{Exported Constants}

\begin{center}
\begin{tabular}{p{4cm} p{4cm} p{4cm} p{4cm}}
\hline
\textbf{Name} & \textbf{In} & \textbf{Out} & \textbf{Exceptions} \\
\hline


\hline
\end{tabular}
\end{center}

\subsubsection{Exported Access Programs}

\begin{center}
\begin{tabular}{p{4cm} p{4cm} p{4cm} p{4cm}}
\hline
\textbf{Name} & \textbf{In} & \textbf{Out} & \textbf{Exceptions} \\
\hline
restart & - & - & CacheOverloadError \\
to\_tutorial & current\_screen\: $\mathbb{Z}$ & Screen\_Object &  PreviousScreenSaveError, CacheError \\
to\_game & current\_screen\: $\mathbb{Z}$ & Screen\_Object &  PreviousScreenSaveError, CacheError  \\
to\_user\_info & current\_screen\: $\mathbb{Z}$ & Screen\_Object &  PreviousScreenSaveError, CacheError, UserNotLoadedError \\
to\_leaderboard & current\_screen\: $\mathbb{Z}$ & Screen\_Object &  PreviousScreenSaveError, CacheError, DatabaseError \\
quit & - & - & - \\


\hline
\end{tabular}
\end{center}

\newpage

\section{Round Control Module} \mref{mRC} 

\begin{description}
\item[Secrets:] The state and progression of a round and the tasks within it (e.g., the selected cell and the player's answers).
\item[Services:] Plays a round given its plane, orientation, starter, and line types; carries out the player's actions and checks them against the ruleset; provides snapshots to \mref{mGDM} for autosave and resume; and records the result through \mref{mUDM} unless the round is marked as a tutorial round. Its read-only access programs provide the round state to \mref{mRE}.
\end{description}

\subsection{Module}

rc

\subsection{Uses}

 Game Data Management Module \mref{mGDM}

 User Data Management Module \mref{mUDM}

 Rendering Module \mref{mRE}


\subsection{Syntax}

\subsubsection{Exported Constants}

\begin{center}
\begin{tabular}{p{4cm} p{4cm} p{4cm} p{4cm}}
\hline
\textbf{Name} & \textbf{In} & \textbf{Out} & \textbf{Exceptions} \\
\hline


\hline
\end{tabular}
\end{center}

\subsubsection{Exported Access Programs}

\begin{center}
\begin{tabular}{p{4cm} p{4cm} p{4cm} p{4cm}}
\hline
\textbf{Name} & \textbf{In} & \textbf{Out} & \textbf{Exceptions} \\
\hline
start\_round & plane: \{$\mathbb{Z}$, $\mathbb{Z}$\}, orientation: $\mathbb{Z}$, starter: Cube, level: $\mathbb{Z}$ & 3D Coordinates List & MissingInfoError, ConversionError, ListOutOfBoundsError \\
detect\_input & User Input Object & - & StateError, SaveError \\
verify\_action & - & Boolean & StateError \\
appply\_action & User Input Object & Data List, Score List, 3D Coordinates List & DatabaseError, TutorialRoundError, RenderError \\

\hline
\end{tabular}
\end{center}


\newpage

\section{Ruleset Module} \mref{mR} 

\begin{description}
\item[Secrets:] The rules of a gameplay round including the number and sequence of tasks, cube order, number-base, additional number counts for each task, valid line-types (e.g., complement, diagonal), and which scoring scheme applies.
\end{description}

\subsection{Module}

ruleset

\subsection{Uses}


\subsection{Syntax}

\subsubsection{Exported Constants}

\begin{center}
\begin{tabular}{p{10cm} p{4cm}}
\hline
\textbf{Name} & \textbf{Value}\\
COMPLEMENT & 1 \\
HALF\_ORTHOGONAL\_X & 2 \\
HALF\_ORTHOGONAL\_Y & 3 \\
HALF\_ORTHOGONAL\_Z & 4 \\
ORTHOGONAL\_X & 5 \\
ORTHOGONAL\_Y & 6 \\
ORTHOGONAL\_Z & 7 \\
DIAGONAL\_XY\_SAME & 8 \\
DIAGONAL\_XY\_DIFF & 9 \\
DIAGONAL\_XZ\_SAME & 10 \\
DIAGONAL\_XZ\_DIFF & 11 \\
DIAGONAL\_YZ\_SAME & 12 \\
DIAGONAL\_YZ\_DIFF & 13 \\
TRIAGONAL\_XYZ\_SAME & 14 \\
TRIAGONAL\_XY\_SAME\_Z\_DIFF & 15 \\
TRIAGONAL\_XZ\_SAME\_Y\_DIFF & 16 \\
TRIAGONAL\_YZ\_SAME\_X\_DIFF & 17 \\
number\_base&8 \\
tasks\_per\_round&8 \\
complement\_sum&0o1001 \\
half\_line\_sum&0o2002 \\
full\_line\_sum&0o4004 \\
half\_line\_extra\_counts&range(1, 3) \\
full\_line\_extra\_counts&range(1, 6) \\
\hline


\hline
\end{tabular}
\end{center}

\subsubsection{Exported Access Programs}

\begin{center}
\begin{tabular}{p{4cm} p{4cm} p{4cm} p{4cm}}
\hline
\textbf{Name} & \textbf{In} & \textbf{Out} & \textbf{Exceptions} \\
\hline
line\_sum & task\_type: TaskType & magic \_ sum: $\mathbb{Z}$ & MissingTypeError, DatabaseMissingCubeTypeError \\
has\_addition\_step & task\_type: TaskType & Boolean & MissingTypeError, DatabaseMissingStepError \\
allowed\_extra\_count & task\_type: TaskType & extra\_nums: List & MissingTypeError, DatabaseMissingListError, OutOfBoundsExtraNumsError \\
\hline
\end{tabular}
\end{center}


\newpage

\section{Scoring Module} \mref{mSC} 

\begin{description}
\item[Secrets:] The available scoring schemes, and the algorithms each uses to calculate task/round scores and performance metrics.
\end{description}

\subsection{Module}

score

\subsection{Uses}


\subsection{Syntax}

\subsubsection{Exported Constants}

\begin{center}
\begin{tabular}{p{10cm} p{4cm}}
\hline
\textbf{Name} & \textbf{Value}\\
DECIMAL\_SCORE\_ADDITION & 10 \\
OCTAL\_SCORE\_ADDITION & 14 \\
DECIMAL\_TIME\_MULT & 1350 \\
OCTAL\_TIME\_MULT & 2506 \\
\hline


\hline
\end{tabular}
\end{center}

\subsubsection{Exported Access Programs}

\begin{center}
\begin{tabular}{p{4cm} p{4cm} p{4cm} p{4cm}}
\hline
\textbf{Name} & \textbf{In} & \textbf{Out} & \textbf{Exceptions} \\
\hline
spawn\_stopwatch & - & - & APINotAvailableError, StopwatchAlreadySpawnedError \\
get\_time & - & time: $\mathbb{R}$ & APINotAvailableError, RestartError \\
add\_to\_score & task\_base: char & - & AdditionError, PreviousAdditionBaseDifferentError, TimeError \\
get\_final\_score & - & final\_score: $\mathbb{N}$ & TimeError, MissingScoringModuleError, DifferentBaseError \\
\hline
\end{tabular}
\end{center}

\newpage

\section{User Management Module} \mref{mUM} 

\begin{description}
\item[Secrets:] The algorithms used to authenticate users and manage user accounts (e.g., create/delete account, login/logout, password reset), to keep track of the current user's session (logged in or playing as a guest), and to compute a user's personal statistics (e.g., completed rounds, highest and average score, fastest round time) from their recorded round results.
\item[Services:] Carries out account actions, and on login or logout sets the current user (or guest) in \mref{mUDM}. Computes personal statistics from the results read from \mref{mUDM} when they are viewed.
\end{description}

\subsection{Module}

user\_management

\subsection{Uses}

User Data Management Module \mref{mUDM}

\subsection{Syntax}

\subsubsection{Exported Constants}

\begin{center}
\begin{tabular}{p{10cm} p{4cm}}
\hline
\textbf{Name} & \textbf{Value}\\
\hline


\hline
\end{tabular}
\end{center}

\subsubsection{Exported Access Programs}

\begin{center}
\begin{tabular}{p{4cm} p{4cm} p{4cm} p{4cm}}
\hline
\textbf{Name} & \textbf{In} & \textbf{Out} & \textbf{Exceptions} \\
\hline
login & username: string, password: string & user\_info: list & UserDatabaseNotAvailable, IncorrectInformationError \\
logout & - & Boolean & - \\
create\_account & username: string, password: string, confirm\_password: string & Boolean & UsernameTakenError, PasswordDoesntMatchError, UserDatabaseNotAvailable \\
delete\_account & - & - & UserDatabaseNotAvailable \\
password\_reset & password: string, confirm\_password: string & Boolean & PasswordDoesntMatchError, UserDatabaseNotAvailable \\
user\_progress\_cache & task\_information: Object & - & CacheFullError, IncorrectInformationError \\
storage\_conversion & final\_information : cache Object & user\_game\_data: arr & ConversionError, CacheEmptyError, UserDatabaseNotAvailable\\
\hline
\end{tabular}
\end{center}

\newpage

\section{Leaderboard Module} \mref{mL} 

\begin{description}
\item[Secrets:] The data structure and algorithms used to select and rank recorded round results into leaderboards (e.g., by scoring base, score or time, all-time or daily challenge).
\end{description}

\subsection{Module}

user\_management

\subsection{Uses}

User Data Management Module \mref{mUDM}

\subsection{Syntax}

\subsubsection{Exported Constants}

\begin{center}
\begin{tabular}{p{10cm} p{4cm}}
\hline
\textbf{Name} & \textbf{Value}\\
\hline


\hline
\end{tabular}
\end{center}

\subsubsection{Exported Access Programs}

\begin{center}
\begin{tabular}{p{4cm} p{4cm} p{4cm} p{4cm}}
\hline
\textbf{Name} & \textbf{In} & \textbf{Out} & \textbf{Exceptions} \\
\hline
get\_leaderboards & - & leaderboards: list of Leaderboard Objects & DailyChallengeDatabaseNotAvailable \\
get\_leaderboard & leaderboards: list of Leaderboard Objects & Leaderboard Object & SeperationError \\
\hline
\end{tabular}
\end{center}

\newpage

\section{Daily Challenge Module} \mref{mDC} 


\begin{description}
\item[Secrets:] How the day's round is obtained from the online server and played on the user's device.
\item[Services:] Fetches the day's round from \mref{mDCM} through \mref{mSCM} and plays it through \mref{mRC}, marked as a daily challenge round so its result is recorded for that day's leaderboard.
\end{description}

\subsection{Module}

daily\_challenge

\subsection{Uses}

Online Daily Challenge Manager Module (\mref{mDCM})

Server Communication Module (\mref{mSCM})

Round Control Module (\mref{mRC})

\subsection{Syntax}

\subsubsection{Exported Constants}

\begin{center}
\begin{tabular}{p{10cm} p{4cm}}
\hline
\textbf{Name} & \textbf{Value}\\
\hline


\hline
\end{tabular}
\end{center}

\subsubsection{Exported Access Programs}

\begin{center}
\begin{tabular}{p{4cm} p{4cm} p{4cm} p{4cm}}
\hline
\textbf{Name} & \textbf{In} & \textbf{Out} & \textbf{Exceptions} \\
\hline
get\_round & - & plane: \{$\mathbb{Z}$, $\mathbb{Z}$\}, orientation: $\mathbb{Z}$, starter: Cube, level: $\mathbb{Z}$ & DailyChallengeDatabaseNotAvailable, TypeError \\
start\_round & plane: \{$\mathbb{Z}$, $\mathbb{Z}$\}, orientation: $\mathbb{Z}$, starter: Cube, level: $\mathbb{Z}$ & - & GameInformationError \\
save\_round\_information & final\_score: $\mathbb{N}$ & - & GameNotFinishedError, DataNotSavedError, CheatingError \\
\hline
\end{tabular}
\end{center}


\newpage

\section{Tutorial Module} \mref{mT} 


\begin{description}
\item[Secrets:] The content and order of the tutorial's steps, and the fixed round they are played on.
\item[Services:] Plays the fixed round through \mref{mRC}, marked as a tutorial round so its result is not recorded, and provides the instruction for the step the player has reached.
\end{description}

\subsection{Module}

tutorial

\subsection{Uses}

Round Control Module \mref{mRC}

\subsection{Syntax}

\subsubsection{Exported Constants}

\begin{center}
\begin{tabular}{p{10cm} p{4cm}}
\hline
\textbf{Name} & \textbf{Value}\\
\hline


\hline
\end{tabular}
\end{center}

\subsubsection{Exported Access Programs}

\begin{center}
\begin{tabular}{p{4cm} p{4cm} p{4cm} p{4cm}}
\hline
\textbf{Name} & \textbf{In} & \textbf{Out} & \textbf{Exceptions} \\
\hline
start\_magic\_cube & - & Navigation Object, 3D Coordinates List & TutorialNotLoadedError, NavigationUIError, MissingExamplesError \\
start\_octal & - & Navigation Object, 3D Coordinates List & TutorialNotLoadedError, NavigationUIError, MissingExamplesError \\
completion\_check & - & Boolean & UserDatabaseError \\
\hline
\end{tabular}
\end{center}

\newpage

\section{Settings Module} \mref{mSE} 


\begin{description}
\item[Secrets:] The application's settings and the algorithms to access and modify them (e.g., turning off animations and configuring the play-time reminder).
\end{description}

\subsection{Module}

settings

\subsection{Uses}


\subsection{Syntax}

\subsubsection{Exported Constants}

\begin{center}
\begin{tabular}{p{10cm} p{4cm}}
\hline
\textbf{Name} & \textbf{Value}\\
\hline


\hline
\end{tabular}
\end{center}

\subsubsection{Exported Access Programs}

\begin{center}
\begin{tabular}{p{4cm} p{4cm} p{4cm} p{4cm}}
\hline
\textbf{Name} & \textbf{In} & \textbf{Out} & \textbf{Exceptions} \\
\hline
change\_settings & base: char, order: char  & Boolean & NoCubeFoundError \\
\hline
\end{tabular}
\end{center}


\newpage

\section{Cube Representation Module} \mref{mCR} 


\begin{description}
\item[Secrets:] The data structure used to store the numbers of the magic cube and look them up by cell.
\item[Services:] Builds a cube of a given size from its numbers, checks that each number appears exactly once, and returns the number in a given cell. The cube size is passed in by the caller, so this module does not use \mref{mR}.
\end{description}

\subsection{Module}

cube\_rep

\subsection{Uses}


\subsection{Syntax}

\subsubsection{Exported Constants}

\begin{center}
\begin{tabular}{p{10cm} p{4cm}}
\hline
\textbf{Name} & \textbf{Value}\\
\hline


\hline
\end{tabular}
\end{center}

\subsubsection{Exported Access Programs}

\begin{center}
\begin{tabular}{p{4cm} p{4cm} p{4cm} p{4cm}}
\hline
\textbf{Name} & \textbf{In} & \textbf{Out} & \textbf{Exceptions} \\
\hline
init & values: List of List of List of Cell Objects  & - & CubeBoundsError, InitializationError \\
value\_at & cell: Coord Object  & value: $\mathbb{R}$ & CubeBoundsError, EmptyCellError \\
check\_shape & values: List of List of List of Cell Objects & Boolean & IncorrectShapeError\\
check\_numbers & - & Boolean & NumAppearsMoreThanOnceError, NotAllNumbersAppearErrpr\\
\hline
\end{tabular}
\end{center}

\newpage

\section{Game Data Management Module} \mref{mGDM} 


\begin{description}
\item[Secrets:] The algorithms to save, modify, delete, and load snapshots of the in-progress round, stored on the device for autosave and resume.
\end{description}

\subsection{Module}

game\_data\_mgmt

\subsection{Uses}


\subsection{Syntax}

\subsubsection{Exported Constants}

\begin{center}
\begin{tabular}{p{10cm} p{4cm}}
\hline
\textbf{Name} & \textbf{Value}\\
\hline


\hline
\end{tabular}
\end{center}

\subsubsection{Exported Access Programs}

\begin{center}
\begin{tabular}{p{4cm} p{4cm} p{4cm} p{4cm}}
\hline
\textbf{Name} & \textbf{In} & \textbf{Out} & \textbf{Exceptions} \\
\hline
save\_current\_round & game\_information: Round Control Object, plane: \{$\mathbb{Z}$, $\mathbb{Z}$\}, orientation: $\mathbb{Z}$, starter: Cube, level: $\mathbb{Z}$, 3D Coordinates List  & round\_cache: Cache Object & CacheFullError, IncorrectGameError \\
remove\_session & - & - & UnableToDeleteError \\
resume\_game & cache: Cache Object & game\_information: Round Control Object, plane: \{$\mathbb{Z}$, $\mathbb{Z}$\}, orientation: $\mathbb{Z}$, starter: Cube, level: $\mathbb{Z}$, 3D Coordinates List & CacheEmptyError \\
\hline
\end{tabular}
\end{center}


\newpage

\section{Round Generation Module} \mref{mRG} 


\begin{description}
\item[Secrets:] The algorithms used to choose a round's plane, orientation, starter, and the line type for each task, consistent with the ruleset.
\end{description}

\subsection{Module}

round\_generator

\subsection{Uses}


\subsection{Syntax}

\subsubsection{Exported Constants}

\begin{center}
\begin{tabular}{p{10cm} p{4cm}}
\hline
\textbf{Name} & \textbf{Value}\\
\hline


\hline
\end{tabular}
\end{center}

\subsubsection{Exported Access Programs}

\begin{center}
\begin{tabular}{p{4cm} p{4cm} p{4cm} p{4cm}}
\hline
\textbf{Name} & \textbf{In} & \textbf{Out} & \textbf{Exceptions} \\
\hline
new\_round & rules: Ruleset Object, level: $\mathbb{Z}$, 3D Coordinates List  & - & RulesError, IncorrectGameError \\
change\_plane & new\_plane: User Input Object & plane: \{$\mathbb{Z}$, $\mathbb{Z}$\} & IncorrectFormationError \\
change\_orientation & new\_orientation: User Input Object & orientation: $\mathbb{Z}$ & IncorrectFormationError \\
change\_starter & starter\_formation: string & starter: Cube & IncorrectFormationError \\
change\_type & new\_type: char & 3D Coordinates List, rules: Ruleset Object & IncorrectFormationError \\
\hline
\end{tabular}
\end{center}


\newpage

\section{Line Geometry Module} \mref{mTG} 


\begin{description}
\item[Secrets:] The algorithms used to compute the cells of a line in the cube from a starter and a line type, including wrapping for split lines, whether a line is continuous, and the order in which the line is traced.
\item[Services:] Given a cube size, a starter cell, and a line type, returns the line's other cells in the order they are traced, and whether the line is continuous or split. The cube size is passed in by the caller, so this module does not use \mref{mR}.
\end{description}

\subsection{Module}

line 

\subsection{Uses}


\subsection{Syntax}

\subsubsection{Exported Constants}

\begin{center}
\begin{tabular}{p{10cm} p{4cm}}
\hline
\textbf{Name} & \textbf{Value}\\
\hline
complement\_sum&0o1001 \\
half\_line\_sum&0o2002 \\
full\_line\_sum&0o4004 \\
half\_line\_extra\_counts&range(1, 3) \\
full\_line\_extra\_counts&range(1, 6) \\

\hline
\end{tabular}
\end{center}

\subsubsection{Exported Access Programs}

\begin{center}
\begin{tabular}{p{4cm} p{4cm} p{4cm} p{4cm}}
\hline
\textbf{Name} & \textbf{In} & \textbf{Out} & \textbf{Exceptions} \\
\hline
find\_line & level: $\mathbb{Z}$, 3D Coordinates List, magic\_constant: $\mathbb{Z}$, rules: Ruleset Object  & Line: List of Cell Objects & LineCantBeFoundError \\
\hline
\end{tabular}
\end{center}

\newpage

\section{User Data Management Module} \mref{mUDM} 


\begin{description}
\item[Secrets:] The data structure and algorithms to save, modify, delete, and load user data, including account records, settings, used starters, and recorded round results (the source of personal statistics and leaderboard rankings).
\item[Services:] Saves and loads data for the current user set by \mref{mUM}. A guest's data is kept only on the device; a logged-in user's data is kept on the server through \mref{mSCM}. Also loads the recorded round results of all users from the server, for \mref{mL} to rank.
\end{description}

\subsection{Module}

user\_data\_mgmt 

\subsection{Uses}

User Management Module \mref{mUM}

Server Communication Module \mref{mSCM}

Leaderboard Module \mref{mL}

\subsection{Syntax}

\subsubsection{Exported Constants}

\begin{center}
\begin{tabular}{p{10cm} p{4cm}}
\hline
\textbf{Name} & \textbf{Value}\\
\hline

\hline
\end{tabular}
\end{center}

\subsubsection{Exported Access Programs}

\begin{center}
\begin{tabular}{p{4cm} p{4cm} p{4cm} p{4cm}}
\hline
\textbf{Name} & \textbf{In} & \textbf{Out} & \textbf{Exceptions} \\
\hline
save\_data & signed\_in\_user: Boolean  & confirmation: cache Object & UserDatabaseError, CacheFullError \\
load\_data & signed\_in\_user: Boolean, user\_id: String  & previous\_game: cache Object, user\_info: Record & UserDatabaseError, NotSignedInError \\
convert\_all\_data & all\_information: database Object & ranked\_information: List of Record & UserDatabaseError, ConversionError \\
\hline
\end{tabular}
\end{center}

\newpage

\section{Number Base Translation Module} \mref{mNBT} 


\begin{description}
\item[Secrets:] The algorithms to convert between the decimal-base values used in gameplay logic and their external user representation in the given number base.
\end{description}

\subsection{Module}

base\_translation

\subsection{Uses}


\subsection{Syntax}

\subsubsection{Exported Constants}

\begin{center}
\begin{tabular}{p{10cm} p{4cm}}
\hline
\textbf{Name} & \textbf{Value}\\
\hline
DIGITS & "0123456789ABCDEF" \\
\hline
\end{tabular}
\end{center}

\subsubsection{Exported Access Programs}

\begin{center}
\begin{tabular}{p{4cm} p{4cm} p{4cm} p{4cm}}
\hline
\textbf{Name} & \textbf{In} & \textbf{Out} & \textbf{Exceptions} \\
\hline
to\_display & value: $\mathbb{Z}$, base: $\mathbb{Z}$ & string\_representation: String & ConversionError \\
parse\_answer & text: String, base: $\mathbb{Z}$  & number: $\mathbb{Z}$ & ConversionError, EmptyStringError \\
check\_base & base: $\mathbb{Z}$ & Boolean & BaseNotSupportedError \\
\hline
\end{tabular}
\end{center}


\newpage

\section{Server Communication Module} \mref{mSCM} 


\begin{description}
\item[Secrets:] The protocol and interface used to exchange data with the online server, including the server's request formats, secure transport, and the handling of failed requests.
\end{description}

\subsection{Module}

server\_cm

\subsection{Uses}


\subsection{Syntax}

\subsubsection{Exported Constants}

\begin{center}
\begin{tabular}{p{10cm} p{4cm}}
\hline
\textbf{Name} & \textbf{Value}\\
\hline

\hline
\end{tabular}
\end{center}

\subsubsection{Exported Access Programs}

\begin{center}
\begin{tabular}{p{4cm} p{4cm} p{4cm} p{4cm}}
\hline
\textbf{Name} & \textbf{In} & \textbf{Out} & \textbf{Exceptions} \\
\hline
ping\_server & server\_info: Server Object & active: Boolean & DatabaseNotFoundError, DatabaseNotActiveError \\
encrypt\_data & user\_information: Object, key: String & encrypted\_user\_information: bytes & EncryptionError \\
decrypt\_data & encrypted\_user\_information: bytes, key: String & user\_information: Server Object (like JSON) & DecryptionError \\
add\_to\_server & server\_info: Server Object, data: Object & confirmation: Server Object & DatabaseNotFoundError, DatabaseNotActiveError, WrongSchemaError \\ 
read\_db & query: SQL query, server\_info: Server Object  & returned\_data: bytes & DatabaseNotFoundError, DatabaseNotActiveError, WrongQueryError \\
\hline
\end{tabular}
\end{center}


\newpage

\section{Online Daily Challenge Manager Module} \mref{mDCM} 


\begin{description}
\item[Secrets:] The algorithms used to create, schedule, and store the daily round on the online server so that every user receives the same round.
\item[Services:] Creates one round each day using \mref{mRG} and the daily challenge ruleset from \mref{mR}, and serves that round to every client that requests it. Runs on the online server, not on the user's device.
\item[Implemented By:]
\item[Type of Module:]
\end{description}

\subsection{Module}

odcm

\subsection{Uses}

Round Generation Module {\mref{mRG}}
Ruleset Module {\mref{mR}}


\subsection{Syntax}

\subsubsection{Exported Constants}

\begin{center}
\begin{tabular}{p{10cm} p{4cm}}
\hline
\textbf{Name} & \textbf{Value}\\
\hline

\hline
\end{tabular}
\end{center}

\subsubsection{Exported Access Programs}

\begin{center}
\begin{tabular}{p{4cm} p{4cm} p{4cm} p{4cm}}
\hline
\textbf{Name} & \textbf{In} & \textbf{Out} & \textbf{Exceptions} \\
\hline
create\_dc & rules: Ruleset Object & daily\_challenge: DailyChallenge Object, starter: Cube, level: $\mathbb{Z}$, 3D Coordinates List & InvalidRulesError \\
start\_dc & daily\_challenge: DailyChallenge Object, starter: Cube, level: $\mathbb{Z}$, 3D Coordinates List, dc\_server\_connectrion: Server Object & - & DatabaseNotAvailable, CheatingDetectedError \\
\hline
\end{tabular}
\end{center}

\bibliographystyle {plainnat}
\bibliography {../../../refs/References}

\newpage

\newpage{}

\section{Appendix --- Generative AI Use Disclosure}

\wss{ChatGPT (version 5) was used to brainstorm ideas for the template for this disclose section.}

\subsection{AI Tools Used}

\wss{Give the name of the AI tool(s) used and the version.}

\subsection{How and Where Used}

\wss{Explain what the tool(s) were used for.  Examples include: brainstorming,
explaining a technical concept, reviewing the document, generating or modifying
text, generating or modifying diagrams, identifying errors or omissions or
inconsistencies.}

\wss{For each of the uses, identify which parts of the document were affected.}

\wss{You do not need to report every interaction, but you should disclose the 
uses that materially contributed to your work}

\subsection{Verification of AI-Generated Content}

\wss{\begin{itemize}
    \item Describe how the team checked AI-generated or AI-assisted material
    before including it in the document.
    \item In particular, verify factual claims, technical assertions,
    requirements, calculations, references, and other information for which an
    AI system may produce incorrect or fabricated results.
    \item \textbf{The team is responsible for the accuracy, completeness,
consistency, and appropriateness of the submitted document, regardless of
whether material was generated by a human or an AI system.}
\end{itemize}}

\newpage{}

\section{Appendix --- Reflection}

\wss{Not required for CAS 741 projects}

The information in this section will be used to evaluate the team members on the
graduate attribute of Problem Analysis and Design.

\input{../../Reflection.text}

\begin{enumerate}
  \item What went well while writing this deliverable? 
  \item What pain points did you experience during this deliverable, and how
    did you resolve them?
  \item Which of your design decisions stemmed from speaking to your client(s)
  or a proxy (e.g. your peers, stakeholders, potential users)? For those that
  were not, why, and where did they come from?
  \item While creating the design doc, what parts of your other documents (e.g.
  requirements, hazard analysis, etc), it any, needed to be changed, and why?
  \item What are the limitations of your solution?  Put another way, given
  unlimited resources, what could you do to make the project better? (LO\_ProbSolutions)
  \item Give a brief overview of other design solutions you considered.  What
  are the benefits and tradeoffs of those other designs compared with the chosen
  design?  From all the potential options, why did you select the documented design?
  (LO\_Explores)
  \item (After you have implemented another team's module, which means this
  isn't filled in until after the original deadline). What did you learn by
  implementing another team's module? Were all the details you needed in the
  documentation, or did you need to make assumptions, or ask the other team
  questions?  If your team also had another team implement one of your modules,
  what was this experience like?  Are there things in your documentation you
  could have changed to make the process go more smoothly for when an
  ``outsider'' completes some of the implementation?

  Your reflection should be about concepts learned in your SE education, such as modular decomposition, cohesion, coupling, information hiding, etc.  You should critique the quality and utility of the specification you received.

  You should create a PR for your implemented module.  A link to the PR should be part of your reflection.

\end{enumerate}

\end{document}